% !TeX program = pdflatex
% !TeX spellcheck = en_GB



%%%%%%%%%%%%%%%%%%%%%%%%%%%%%%%%%%%%%%%%%
%%  FUNDAMENTAL SETTINGS AND PACKAGES  %%
%%%%%%%%%%%%%%%%%%%%%%%%%%%%%%%%%%%%%%%%%


\documentclass[11pt, a4paper, notopicbreak, titleabove]{simplecv}

\let \oldtopic \topic
\renewcommand{\topic}{%
	\oldtopic%
	\setlength{\parskip}{0.45ex}%
	\setlength{\itemsep}{0.75ex}%
}

\usepackage[american, ngerman]{babel} % added ngerman because \selectlanguage{ngerman} is used later
\usepackage[american, cleanlook]{isodate}


\newcommand{\myname}{Gunilla Kaibel}

\usepackage{xpatch}
\usepackage[utf8]{inputenc}  % so that umlauts can be input without having to use TeX code
\usepackage{ragged2e}
\usepackage{changepage}
\usepackage{multicol}

\usepackage{calc}
\usepackage{tikz}
\usepackage{tcolorbox}
\usetikzlibrary{positioning}
\usepackage{xcolor}
\definecolor{UBonnBlue}   {RGB}{  7,  82, 154}
\definecolor{UBonnYellow} {RGB}{234, 185,  12}
\definecolor{UBonnGray}   {RGB}{144, 144, 133}
\definecolor{darkgray}    {RGB}{102, 102, 102}
\definecolor{darkred}     {RGB}{153,   0,   0}
\definecolor{neutralgreen}{RGB}{ 28, 166,   0}

\definecolor{QueenslandPurple}{HTML}{51247A}
\definecolor{QueenslandGrey}{HTML}{938884}

\colorlet{SpotColor}{UBonnBlue}%{QueenslandPurple}

\usepackage{authoraftertitle}

\usepackage[hyphens]{url}
\usepackage[unicode = true]{hyperref}
\hypersetup{
	colorlinks = true,
	allcolors = black,
	urlcolor = SpotColor,
	bookmarks = true,
	bookmarksnumbered = false,
	bookmarksopenlevel = 1,
	pdfstartview = Fit,
	pdfpagelayout = SinglePage,
	plainpages = false,
	pdfpagelabels = true,
	pdfauthor = \myname,
	pdftitle = {\myname's Curriculum Vitae}
}
\urlstyle{same}
\Urlmuskip = 0mu\relax  % Prevent additional whitespace before/after breakable characters in URLs
\newcommand{\email}[1]{\href{mailto:#1}{\nolinkurl{#1}}}
\newcommand{\doi}[1]{\url{https://doi.org/#1}}




%%%%%%%%%%%%%%%%%%%
%%  PAGE LAYOUT  %%
%%%%%%%%%%%%%%%%%%%


\raggedbottom

\usepackage{geometry}
\geometry{
	a4paper,
	left = 5mm,
	right = 25mm,
	top = 20mm,
	bottom = 20mm
}
\renewcommand{\topicmargin}{42.5mm}
% Make the label right-aligned:
\makeatletter
\renewcommand{\@topic@makelabel}[1]{\hfill\cv@tlab@fnt #1\quad}
\renewcommand{\cv@do@title}{%
	\par\bigskip
	{\noindent\hspace{\topicmargin}\cv@tit@fnt\cv@tit}%
}
\makeatother

% Slightly increase the line spacing
\renewcommand{\baselinestretch}{1.03}

% Title spacing
\usepackage{titlesec}
\titlespacing*{\section}{\topicmargin}{0.65cm}{0.35cm}

% Footer ==>
\usepackage{pageslts}
\pagenumbering{arabic}
\usepackage{fancyhdr}
\AtBeginDocument{%
	\title{\MyAuthor}%
}
\renewcommand{\headrulewidth}{0pt}
\pagestyle{fancy}
\fancyhf{}
\lfoot{%
	\footnotesize\sffamily\color{SpotColor}\\[-0.15cm]%
	\makebox[\topicmargin][r]{Seite~\thepage\ von~\pageref*{LastPage}\enspace\quad}%
	\textbf{\MyTitle}\enspace\quad%
	\MyDate%
}
% <==

\newenvironment{preventpagebreak}
	{\par\begin{minipage}[t]{\dimexpr\textwidth - \topicmargin\relax}}
	{\end{minipage}\par\vspace{\parskip}}

% Allow for multiline topic labels:
\newcommand{\multiline}[1]{%
	\parbox[t][10pt][t]{\topicmargin-50pt}{\RaggedLeft%
		#1%
		\normalsize\par%
			% So that \baselineskip is of the normal size, see https://tex.stackexchange.com/a/408544
	}%
}


%%%%%%%%%%%%%%%%%%%%%
%%  FONT SETTINGS  %%
%%%%%%%%%%%%%%%%%%%%%


\iftutex
	\usepackage[protrusion=true, expansion=false]{microtype}
	\usepackage{fontspec}
	%\setmainfont[
	%	BoldFont = Bookerly-Bold,
	%	BoldItalicFont = Bookerly-BoldItalic,
	%	ItalicFont = Bookerly-Italic,
	%	Extension = .ttf,
	%	Path = {./Amazon_Typefaces_Complete_Font_Set_Mar2020/Bookerly/},
	%	Scale = 0.875,
	%]{Bookerly-Regular}
	%\setsansfont[
	%	BoldFont = AmazonEmber_Bd,
	%	BoldItalicFont = AmazonEmber_BdIt,
	%	ItalicFont = AmazonEmber_RgIt,
	%	Extension = .ttf,
	%	Path = {./Amazon_Typefaces_Complete_Font_Set_Mar2020/Ember/},
	%	Scale = 0.875,
	%]{AmazonEmber_Rg}
	\setmainfont[
		BoldFont = MuseoSans_700,
		BoldItalicFont = MuseoSans_700_Italic,
		ItalicFont = MuseoSans_300_Italic,
		Extension = .otf,
		Numbers = Tabular,
		Path = ./MuseoSans/,
		Scale = 0.94,
	]{MuseoSans_300}
	\setsansfont[
		BoldFont = MuseoSans_700,
		BoldItalicFont = MuseoSans_700_Italic,
		ItalicFont = MuseoSans_300_Italic,
		Extension = .otf,
		Numbers = Tabular,
		Path = ./MuseoSans/,
		Scale = 0.94,
	]{MuseoSans_300}
	\topiclabelfont{%
		\small\sffamily\color{SpotColor}%
	}
\else
	\usepackage[protrusion=true, expansion=false, kerning=true]{microtype}
	\DisableLigatures[f]{family = {rm*}}
	% Disable the f* ligatures for Charter because the font does not provide sufficient support.
	\DisableLigatures[f]{family = {tt*}}
	% Disable the f* ligatures for Fira Mono because they make no sense for a monospaced font.
	\SetExtraKerning[unit=space]
		{encoding=*, family=*, series=*, size={*, normalsize, footnotesize}, font = */*/*/*/*}
		{\textemdash = {325, 325},
		 / = {100, 100},
		 : = { 50,   0},
		 ; = { 50,   0}}
	% \input{Preamble_mweights}
	%\usepackage[scale=0.88]{plex-sans}
	%\usepackage[semibold, scale=0.88]{plex-serif}
	%\topiclabelfont{\sffamily\renewcommand{\bfseries}{\fontseries{semibold}\selectfont}}
	%\renewcommand{\bfseries}{\fontseries{semibold}\selectfont}
	% \usepackage[charter, greeklowercase=italicized, greekuppercase=italicized, sfscaled=false, ttscaled=false]{mathdesign}
	%\usepackage[scaled=.96, lining, sups]{XCharter}  % option ``sups'' not working properly with subimport
	% NOTE (2025-10-01): Removed newpx (Palatino) math/text package because it clashes with mathdesign
	% Both mathdesign and newpx define a script alphabet (\mathscr), triggering:
	%   ! LaTeX Error: Command \mathscr already defined.
	% If you prefer Palatino (newpx) over Charter (mathdesign), instead comment out mathdesign
	% above and uncomment the next line. Do NOT load both simultaneously.
	\usepackage[scaled=0.95, lining]{newpx}
	%\usepackage[book, semibold, lining, tabular, scale=0.92]{FiraSans}
	\usepackage[scaled=.85, medium, tabular]{montserrat}
	\renewcommand{\rmfamily}{\sffamily}
	\topiclabelfont{%
		\small\sffamily%
		%\renewcommand{\bfseries}{\fontseries{medium}\selectfont}
		\color{SpotColor}%
	}
\fi

% Enabling slightly reduced font for CAPS:
\usepackage{relsize}
\renewcommand\RSpercentTolerance{1}
	\newcommand{\caps}[1]{\textscale{0.96}{\textls[35]{\MakeUppercase{#1}}}}

\titlefont{\Large\sffamily\bfseries\color{SpotColor}}
\newcommand{\mysectionfont}{\large\sffamily\bfseries}
\sectionfont{\mysectionfont\color{SpotColor}}
\topictitlefont{\rmfamily}

\newcommand{\commabbrev}[1]{\textsf{#1}}

\newcommand{\highlight}[1]{%
	\tikz[overlay]
	\node[fill=SpotColor, inner xsep=4pt, text height=0.85*height("Ä1g"), text depth=0.5*depth("Ä1g"), anchor=text, rectangle]{%
		\textcolor{white}{\textbf{#1}}%
		\hspace{-0.33pt}
	};%
	\phantom{\textbf{#1}}%
}

\newlength{\testwidth}
\newcommand{\highlightsection}[1]{%
	\setlength{\testwidth}{\widthof{\mysectionfont #1}}
	\addtolength{\testwidth}{12pt}
	\medskip
	\begin{tcolorbox}[
		colback=SpotColor, colframe=SpotColor,
		width=1\testwidth,
		arc=0pt, outer arc=0pt, boxrule=0pt,
		left=6pt, right=6pt, top=4pt, bottom=4pt, boxsep=0pt
	]%
		\section{\textcolor{white}{\mbox{#1}}}%
	\end{tcolorbox}%
}

\usepackage[fixed]{fontawesome5}
\newcommand{\icon}[1]{%
	\footnotesize\textcolor{SpotColor}{\faIcon[regular]{#1}}
}

\frenchspacing
\sloppy




%%%%%%%%%%%%
%%  BODY  %%
%%%%%%%%%%%%


\begin{document}

\renewcommand{\rmfamily}{\sffamily}

\RaggedRight
%\hfill%
%\raisebox{\baselineskip}{\makebox(0, 0)[rt]{
%	%\includegraphics[height=40mm, width=35mm, trim = 50mm 80mm 50mm 30mm, clip]{me.jpg} %trim={left bottom right top}
%	\includegraphics[height=36mm, %trim = 50mm 80mm 50mm 30mm,
%	clip]{example-grid-100x100pt} %trim={left bottom right top}
%}}%
%\vspace{-\baselinestretch\baselineskip}%
%\vspace{-\baselinestretch\baselineskip}%
%\vspace{-\parskip}%

\author{\myname}
\date{\today}
\maketitle

\setlength{\leftskip}{\topicmargin}
\vspace{-\medskipamount}

%\sffamily
\begingroup
	\color{QueenslandGrey}
	\noindent Doktorandin \\
	\textbf{\color{SpotColor}Deutsches Seminar, Universität Zürich}
\endgroup

\smallskip

\begin{topic}

	\pdfbookmark[1]{Contact information}{contactinfo}
	\item[\multiline{%
		\icon{map} {\normalsize\\}% \icon{envelope} {\normalsize\\}
		\icon{comments} {\normalsize\\}%
		\icon{address-card}%
	}]
	Schönberggasse 2, 8001 Zürich, Schweiz \\
	\email{Gunilla.Kaibel@ds.uzh.ch} \\
	\href{https://www.ds.uzh.ch/apps/p/kaibel}{www.ds.uzh.ch/apps/p/kaibel}

\end{topic}




\section{Akademische Positionen}

\begin{topic}
	\item[2023-- heute]
		Doktorandin {\color{SpotColor}|} Universität Zürich  \par
		\begingroup
			\color{QueenslandGrey}\small%
			\fontseries{regular}\selectfont%
			SNF-/DFG-Projekt 
„Zwischen Erwartungshaltung und Empathie: Expertise-Aushandlung und Verständigungspraktiken in der Online-Wissenschaftskommunikation“ (Prof. Noah Bubenhofer, Prof. Nina Janich, Dr. Michael Bender)\\  
		\endgroup
\end{topic}

\begin{topic}
	\item[2022-- 2023]
		Studentische Mitarbeiterin {\color{SpotColor}|} Universität Marburg \par
		\begingroup
			\color{QueenslandGrey}\small%
			\fontseries{regular}\selectfont%
			Abteilung Semantik des Deutschen Sprachatlas beim Hessen-Nassauischen Wörterbuch in der Lexikographie (Prof. Alfred Lameli, Dr. Bernd Vielsmeier)\\  
		\endgroup
\end{topic}

\begin{topic}
	\item[2022]
		Externe Beraterin {\color{SpotColor}|} Leibniz-Institut für Deutsche Sprache \par
		\begingroup
			\color{QueenslandGrey}\small%
			\fontseries{regular}\selectfont%
			Digitalisierung des Deutschen Fremdwörterbuchs\\  
		\endgroup
\end{topic}

\begin{topic}
	\item[2022]
		Praktikum {\color{SpotColor}|} Leibniz-Institut für Deutsche Sprache \par
		\begingroup
			\color{QueenslandGrey}\small%
			\fontseries{regular}\selectfont%
			im Projekt „Sprache(n) im öffentlichen Raum“ (Prof. Albrecht Plewnia) und in der Lexik (Prof. Stefan Engelberg)\\  
		\endgroup
\end{topic}

\begin{topic}
	\item[2020]
		Praktikum {\color{SpotColor}|} Rheinische Friedrich-Wilhelms-Universität Bonn \par
		\begingroup
			\color{QueenslandGrey}\small%
			\fontseries{regular}\selectfont%
			beim Dialektatlas Mittleres Westdeutschland (Prof. Claudia Wich-Reif)\\  
		\endgroup
\end{topic}

\begin{topic}
	\item[2020]
		Praktikum {\color{SpotColor}|} Adalbert-Stifter-Institut in Linz (AT)  \par
		\begingroup
			\color{QueenslandGrey}\small%
			\fontseries{regular}\selectfont%
			im Bereich Sprachforschung\\  
		\endgroup
\end{topic}

\begin{topic}
	\item[2018--2019]
		Studentische Hilfskraft {\color{SpotColor}|} Rheinische Friedrich-Wilhelms-Universität Bonn  \par
		\begingroup
			\color{QueenslandGrey}\small%
			\fontseries{regular}\selectfont%
			SFB 1167 „Macht und Herrschaft“ im mediävistischen Teilprojekt „Publizistische Zeitklagen“ (Prof. Karina Kellermann)\\  
		\endgroup
\end{topic}


\section{Ausbildung}

\begin{topic}

	\item[2024--2025]
		Hochschuldidaktische Weiterbildung \textit{Teaching Skills} {\color{SpotColor}|} Universität Zürich \par
		\begingroup
			\color{QueenslandGrey}\small%
			%\fontseries{regular}\selectfont%
			% \textit{Dissertation:} ``TBA'' \\  %{\color{SpotColor}|}
		
		\endgroup


	\item[2021--2023]
		Master Linguistik: Kognition und Kommunikation  {\color{SpotColor}|} Philipps-Universität Marburg
		\begingroup
			\color{QueenslandGrey}\small%
			\fontseries{regular}\selectfont%
			Masterarbeit: Dialekt im Diskurs. Eine diachrone Quellenanalyse \\
			Abschlussnote 0,9 \\  
		\endgroup
	\item[2020]
		Erasmus Semester {\color{SpotColor}|} Université de Picardie Jules Verne, Amiens		
	
	\item[2019]
		Sommerakademie Mittelniederdeutsch: Literatur – Sprache – Medien {\color{SpotColor}|} Universität Rostock
	\item[2017--2021]
		Zwei-Fach-Bachelor Germanistik und Kunstgeschichte {\color{SpotColor}|} Rheinische Friedrich-Wilhelms-Universität Bonn \\
		\begingroup
			\color{QueenslandGrey}\small%
			\fontseries{regular}\selectfont%
			Bachelorarbeit: Dialekt in interpersonaler Onlinekommunikation \\ 
			Abschlussnote 1,1 \\ 
		\endgroup
	\item[2008-2016]
		Abitur {\color{SpotColor}|} Viktoriaschule Aachen \\
		\begingroup
			\color{QueenslandGrey}\small%
			\fontseries{regular}\selectfont%
			Abschlussnote 1,3 \\  
		\endgroup
\end{topic}


\section{Forschungsaufenthalte}

\begin{topic}
	\item[2025]
		Forschungsaufenthalt {\color{SpotColor}|} Humboldt-Universität Berlin \par
		\begingroup
			\color{QueenslandGrey}\small%
			\fontseries{regular}\selectfont%
			Gastgeberin Prof. Anke Lüdeling \\  
		\endgroup
	\item[2025]
		Forschungsaufenthalt {\color{SpotColor}|} Universität Uppsala \par
		\begingroup
			\color{QueenslandGrey}\small%
			\fontseries{regular}\selectfont%
			Gastgeber Prof. Michael Prinz \\  
		\endgroup

	

\end{topic}



\section{Publikationen}

\begin{topic}

	\item[Aufsatz]
    Kaibel, Gunilla (2025): Constructing Expertise Through Questions? -- Potentials of Corpus Pragmatic Regression Analyses. In: \textit{Corpus Pragmatics} 10(1), \url{https://doi.org/10.1007/s41701-025-00206-4}

	\item[Aufsatz]
    Knuchel, Daniel / Bojarski, Xenia / Ventre, Davide / Huber, Sonja / Kaibel, Gunilla (2026): Diskurswandel korpuspragmatisch aufspüren. Analyse- und Visualisierungsmethoden im Kontext der data philology. In: Brunner, Annelen / Hansen, Sandra / Lang, Christian / Tu, Ngoc Duyen Tanja / Wolfer, Sascha (Hg.): \textit{Deutsch im Wandel - Tools, Ressourcen und Methoden. Beiträge zur Methodenmesse der 61. Jahrestagung des Leibniz-Instituts für Deutsche Sprache}. Mannheim: IDS-Verlag. Im Druck.

	\item[In Begutachtung]
	\textit{Mal ne ernstgemeinte Frage:} Kann man Verstehen steuern und damit polarisierte Online-Diskussionen deeskalieren?

	\item[In Begutachtung]
	\textit{kein hass sondern ernstgemeinte frage}. Formen und Funktionen metapragmatischer Verstehenslenkung in Online-Diskussionen

	\item[In Begutachtung]
	Mediale Diskurse über Dialekt im Nationalsozialismus. In: \textit{RegioLingua}

\end{topic}



\section{Vorträge}

\begin{topic}

	\item[2026]
	DGPuK-Fachgruppentagung Mediensprache – Mediendiskurse (Landau), KoKoKom-Tagung (Tübingen)

	\item[2025]
	Germanistentag (Braunschweig), diverse Kolloquien (Zürich, Uppsala, Stockholm, Berlin), Posterpräsentation Methodenmesse der IDS-Jahrestagung (Mannheim)

	\item[2024]
	Forum Sprachvariation (Kiel), Posterpräsentation Sprache und Wissen (Heidelberg)

	\item[2021]
	Forum Linguistik (Bonn)

\end{topic}



\section{Lehre}

\begin{topic}

	\item[Dozentin]
	Modul \textit{Bestehende Korpusressourcen} {\color{SpotColor}|} Zürcher Sommerschule für Korpuslinguistik und Korpuspragmatik (ZuKoKo), Universität Zürich {\color{SpotColor}|} 2025
	\par
	\begingroup
		\color{QueenslandGrey}\small%
		Konzeption des Moduls \par
	\endgroup

	\item[Dozentin]
	Bachelorkurs \textit{Einführung in die germanistische Linguistik A} {\color{SpotColor}|} Deutsches Seminar, Universität Zürich {\color{SpotColor}|} Herbstsemester 2024 \par
	\begingroup
		\color{QueenslandGrey}\small%
		Sechs Lektionen komplett selbst konzipiert, eine weitere Lektion für alle Einführungskurse neu konzipiert \par
	\endgroup

\end{topic}

\section{Auszeichnungen und Stipendien}
\begin{topic}

\item[2022--2023]
Hausstipendium der Hessischen Stipendiat*innenanstalt Collegium Philippinum

\item[2019--2021]
Deutschlandstipendium der Universität Bonn

\item[2019--2021]
Erfolgreiche Teilnahme am Honors Program der Universität Bonn

\end{topic}


\section{Eingeworbene Drittmittel}
\begin{topic}
	\item[2025]
	3.130 CHF Reisestipendium {\color{SpotColor}|} Graduate Campus der Universität Zürich  

\end{topic}


\section{Ehrenamtliches Engagement}
\begin{topic}
	\item[2024--heute]
	Mitglied in der Leitung der Konferenz des Wissenschaftlichen Nachwuchses {\color{SpotColor}|} Deutsches Seminar, Universität Zürich

	\item[2021--2023]
	Mitarbeit in verschiedenen Ämtern, u. a. als studentische Vertreterin der Verwaltungskommission {\color{SpotColor}|} Hessische Stipendiat*innenanstalt Collegium Philippinum, Marburg

	\item[2019, 2021]
	Hausleitung {\color{SpotColor}|} Ev.-Theol. Studienhaus Adolf Clarenbach, Bonn

\end{topic}




\section{Kenntnisse und Fähigkeiten}

\begin{topic}

	\item[Fremdsprachen]
	Englisch (C1), Französisch (B2), Schwedisch (A2)

	\item [IT-Kenntnisse]
	Programmiersprachen (python, R), Microsoft Office (Word, Excel, Access), LaTeX, Bild- und Videobearbeitung (GIMP, Camtasia), Suchsprachen (ANNIS Query Language), XML (Oxygen) 


\end{topic}




\end{document}
